\rok{Јануар 2024.}{А}

Други практични тест из Алгоритама и структура података

Други практични тест одржан 09.01.2024.

\zadatak{1}{Конверзија листе у матрицу суседства}
Написати алгоритам који ће вршити директну конверзију листе суседства у матрицу суседства. Алгоритам као улазни параметар треба да прихвати динамички низ чворова, а треба да резултује креираном матрицом.

\textbf{Додатне напомене за решавање задатка:}
\begin{itemize}
	\item У template-у је закоментарисана метода \texttt{ConvertAdjacencyListToAdjacencyMatrix} која треба да садржи имплементацију конверзије листе у матрицу суседства.
	\item У Main методи обезбедити начин за тестирање и проверу нове функционалности.
\end{itemize}



\zadatak{2}{Проналажење минималног индекса понављајуће вредности у низу}
Написати алгоритам временске комплексности $O(n)$ који ће претражити низ целих бројева и пронаћи индекс прве понављајуће вредности у низу.

\textbf{Додатне напомене за решавање задатка:}
\begin{itemize}
	\item Обавезно је користити Hash табелу која је дата у шаблону.
	\item Имплементацију писати у оквиру закоментарисане методе:
	\begin{Verbatim}[fontsize=\footnotesize, numbers=left, xleftmargin=10mm]
public static int findMinIndexOfRepeatingNumber(int[] arr)
	\end{Verbatim}
\end{itemize}

\textbf{Примери:}
\begin{itemize}
	\item \textbf{Улаз 1:} \texttt{[5, 6, 3, 4, 3, 6, 4]}
	\item \textbf{Излаз 1:} 1
	\item \textbf{Улаз 2:} \texttt{[1, 2, 3, 4, 5, 6]}
	\item \textbf{Излаз 2:} -1
\end{itemize}



\zadatak{3}{Проналажење најдуже путање до неког листа стабла}
У оквиру стандардне структуре класе \texttt{BinarySearchTree} написати имплементацију методе која ће омогућити проналажење најдуже путање до неког листа у оквиру стабла.

\textbf{Додатне напомене за решавање задатка:}
\begin{itemize}
	\item У оквиру класе \texttt{BinarySearchTree} креирати јавну методу са потписом:
	\begin{Verbatim}[fontsize=\footnotesize, numbers=left, xleftmargin=10mm]
public void findLongestLeafPath()
	\end{Verbatim}
	\item Из претходно креиране јавне методе треба да се позива приватна метода са потписом:
	\begin{Verbatim}[fontsize=\footnotesize, numbers=left, xleftmargin=10mm]
private void findLongestLeafPathBST(Node curr)
	\end{Verbatim}
	\item Приватна метода треба да садржи имплементацију алгоритма којим ће се омогућити пролазак кроз стабло од корена ка листовима и исписивање чворова који се налазе на најдужој путањи до листа.
	\item Лист представља сваки чвор који нема ниједно дете у структури BST стабла.
	\item Корен стабла се налази на нултом нивоу.
	\item Уколико постоји више путања исте дужине, одабрати „скроз десну“ путању у стаблу.
	\item Алгоритам \textbf{ОБАВЕЗНО} писати у оквиру класе \texttt{BinarySearchTree}.
\end{itemize}

\textbf{Примери:}

\begin{center}
\begin{minipage}{\textwidth}
\centering
\begin{forest}
for tree={
	circle,
	draw,
	thick,
	fill=gray!20,
	minimum size=8mm,
	inner sep=1pt,
	s sep=8mm,
	l sep=12mm,
	edge={-, thick},
	font=\small
}
[8
	[4
		[2
			[, phantom]
			[3]
		]
		[5]
	]
	[9
		[, phantom]
		[11]
	]
]
\end{forest}

\vspace{0.3cm}
\textbf{Слика 1.} Пример BST-а.
\end{minipage}
\end{center}

\begin{itemize}
	\item \textbf{Улаз 1:} BST са слике 1
	\item \textbf{Излаз 1:} \texttt{8 -> 4 -> 2 -> 3}
\end{itemize}

\begin{center}
\begin{minipage}{\textwidth}
\centering
\begin{forest}
for tree={
	circle,
	draw,
	thick,
	fill=gray!20,
	minimum size=8mm,
	inner sep=1pt,
	s sep=8mm,
	l sep=12mm,
	edge={-, thick},
	font=\small
}
[10
	[4
		[2
			[3]
			[, phantom]
		]
		[6
			[5]
			[9]
		]
	]
	[20
		[12
			[, phantom]
			[14]
		]
		[29
			[27]
			[, phantom]
		]
	]
]
\end{forest}

\vspace{0.3cm}
\textbf{Слика 2.} Пример BST-а.
\end{minipage}
\end{center}

\begin{itemize}
	\item \textbf{Улаз 2:} BST са слике 2
	\item \textbf{Излаз 2:} \texttt{10 -> 20 -> 29 -> 27}
\end{itemize}



\sablon{Шаблон другог теста јануар 2024. група А}{2023/Januar/GrupaA/AISP2023_Test2_GrupaA.linq}
